
%%%%%%%%%%%%%%%%%%%%%%%%%%%%%%%%%%%%%%%%%%%%%%%%%%%%%%%%%%%%%
\section{符号说明}

本文主要符号及其含义如下：

\begin{longtable}[]{@{}lll@{}}
\toprule
\textbf{符号} & \textbf{含义} & \textbf{单位/说明} \\
\midrule
\endhead

$n$ & 患者总数 & 个 \\

$m$ & 健康指标维度数 & 个 \\

$X$ & 患者特征矩阵 & $n \times m$ 维 \\

$x_{ij}$ & 第$i$个患者的第$j$个指标值 & 根据指标类型 \\

$Y$ & 目标变量向量 & $n$ 维 \\

$y_i$ & 第$i$个患者的健康评分或风险等级 & 0-100或1-4 \\

$w$ & 模型参数向量 & $m$ 维 \\

$H_i$ & 患者$i$的健康评分 & 0-100分 \\

$R_i$ & 患者$i$的风险等级 & 1(低)-4(极高) \\

$P_i$ & 患者$i$未来1年疾病急性发作概率 & 0-1 \\

$S_j$ & 第$j$个管理方案 & 方案集合 \\

$d(i,j)$ & 患者$i$与$j$的特征相似度 & 0-1 \\

$ACC$ & 模型准确率 & \% \\

$TPR$ & 真阳性率（灵敏度） & \% \\

$FPR$ & 假阳性率 & \% \\

$AUC$ & ROC曲线下面积 & 0-1 \\

$RMSE$ & 均方根误差 & 根据指标 \\

$p$值 & 统计检验的显著性 & < 0.05为显著 \\

\bottomrule
\end{longtable}

本文使用的主要符号及其含义如表\ref{tab:符号说明}所示。

% longtable：可跨页，>{\centering\arraybackslash}p{} 居中，比例按内容调整
\begin{longtable}{>{\centering\arraybackslash}p{0.18\textwidth}>{\centering\arraybackslash}p{0.78\textwidth}}
  \caption{符号说明}
  \label{tab:符号说明} \\
  \toprule
  符号 & 说明 \\
  \midrule
  \endfirsthead
  \caption{符号说明（续）} \\
  \toprule
  符号 & 说明 \\
  \midrule
  \endhead
  \bottomrule
  \endfoot
  ... & ... \\
  ... & ... \\
  ... & ...
\end{longtable}
